\documentclass[../main.tex]{subfiles}
\begin{document}
\setlist[itemize]{topsep=2pt,itemsep=1pt,parsep=0pt,partopsep=0pt}
\setlist[itemize,2]{topsep=1pt,itemsep=0pt,parsep=0pt,partopsep=0pt}
\setlength{\parskip}{3pt}

\chapter{Week 5 — Data Management and Project Planning}
\label{ch:week5}

\textit{Snapshot and part summaries: see the overview on page~\pageref{ov:week5}.}

\section{Part 1 — The Data Management Plan}
\label{sec:dmp}

\begin{itemize}
  \item The \keyterm{DMP} must be written before data collection — organising data only at the end is close to impossible.
  \item \textbf{Section I — Data Description}: data type, file formats, collection method with source and terms of use, processing purpose, storage location, access rights.
  \item \textbf{Section II — Documentation and Data Quality}: collection methodology, electronic lab notebook, README file, data dictionary of variables.
  \item \textbf{Section III — Storage and Backup}: \keyterm{3-2-1 rule} — 3 copies, on 2 storage media, 1 copy in a different physical location. Design this in from the start.
  \item \textbf{Section IV — Legal and Ethical Requirements}: confidential/classified data; human subjects or third-party datasets.
  \item \textbf{Section V — Data Sharing and Long-term Preservation}: what is shared, how (e.g.\ 4TU.ResearchData), and under which licence (CC0, CC-BY, CC BY-SA, CC BY-NC).
  \item \textbf{Section VI — Data Management Responsibilities}: confirm TU Delft leads the project and name who becomes responsible if you leave.
  \item Further rules: descriptive file/variable names, consistent filing system, uncorrected original, open formats, documented workflow.
\end{itemize}

\section{Part 2 — Project Management in Research}
\label{sec:project-mgmt}

\begin{itemize}
  \item Project management delivers organised progress, resource optimisation, early risk mitigation, and effective supervisor communication.
  \item \keyterm{Task estimation}: outline every task, estimate each using comparative, parametric, or bottom-up estimating.
  \item \keyterm{Work Breakdown Structure (WBS)}: decomposes the project into work packages and tasks to the level at which each task can be estimated; estimates sum to the project total.
  \item \keyterm{Resource management}: resources are people (reusable), equipment, and facilities (consumable). Describe required resources and their availability per task.
  \item \keyterm{Project planning}: time-based schedule — only possible once conceptual and technical designs are complete.
  \item \keyterm{Gantt chart}: standard schedule representation; reveals the \keyterm{critical path} — the longest dependent chain that determines the minimum project duration.
\end{itemize}

\section{Reading Material}
\label{sec:reading5}

\begin{partbox}
The Week~5 reading, \textit{Sustainable Data Management} (Pommerening), supplies the principles behind the DMP, centred on the FAIR framework and a classification of research data.
\end{partbox}

\begin{itemize}
  \item \keyterm{FAIR}: findability, accessibility, interoperability, reuse — reuse is the ultimate objective that documentation serves.
  \item Five data categories: observational (irreplaceable, time- and place-bound), experimental (reproducible but costly), simulation, derived, and metadata.
  \item Data outlives the project — sharing is now part of the scientific process, not an afterthought.
  \item Spreadsheet-driven analysis is hard to reconstruct; scripted workflows are easier to repeat and revise.
  \item Missing metadata and workflow documentation eliminates reproducibility, transparency, and reuse.
\end{itemize}

\end{document}
